% Matthew: maximin fairness under group-dependent acceptance.
% \input from a1.tex after sections/setup (shared model, B_c, E_c, r_c, ceiling model, Claim 1).

\subsection{Formulation}
The maximin objective and the lexicographic problem below are taken from FWRRS \cite{fwrrs}. Our change is to the diffusion model (the acceptance probabilities $a_v$), not to the objective. The maximin value is the worst infected community's rate,
\begin{equation}
\mathbb{M}(S_N,S_P) = \min_{c \in \mathcal{C}_{\mathrm{infected}}} r_c(S_P),
\end{equation}
and the fairness-aware blocking problem maximizes total blocking among the maximin-optimal seed sets:
% ---------- Fairness-aware blocking problem ----------
\begin{equation}
\begin{aligned}
\max_{S_P \in \mathcal{S}}\quad & B(S_P) \\
\text{s.t.}\quad & S_P \in \operatorname*{arg\,max}_{S'_P \in \mathcal{S}}\ \mathbb{M}(S_N,S'_P)
\end{aligned}
\end{equation}
\begin{itemize}
\item \textbf{Relation to DP.} Maximin and DP rank seed sets by the same per-community rates and differ only in what they do with them: DP shrinks the spread of the rates, maximin lifts the lowest.
\item \textbf{Two outcomes.} Like DP, $\mathbb{M}$ can be computed on either outcome; we write $\mathbb{M}^{\mathrm{sav}}$ and $\mathbb{M}^{\mathrm{conv}}$.
\end{itemize}

\subsection{Properties and algorithm}
\begin{itemize}
\item \textbf{Monotone.} Fix a realization. Adding a truth seed can only make the truth arrive earlier and the rumor later or not at all (nodes that reject the truth stay inactive until the rumor reaches them), so the truth-active set can only grow and the rumor-active set can only shrink. So both the saved and converted counts in each $c$, and hence $\mathbb{M}$, are monotone in $S_P$.
\item \textbf{Not submodular.} With two infected communities that each need their own seed, one seed leaves $\mathbb{M}=0$ and the second lifts it above $0$, so the second seed gains more than the first. Greedy on $\mathbb{M}$ therefore sees zero gain from every single seed while two or more communities sit at $0$ and no single seed reaches all of them, and has no $1-1/e$ guarantee. Rejection also breaks submodularity of $B$ itself, even with equal acceptance across groups (direction 2).
\item \textbf{Worst-first greedy.} Following FWRRS, at each step we add the candidate with the largest gain in the currently worst community's rate, breaking ties by the gain in $B$. This is a greedy heuristic for the lexicographic objective above, not an exact solution. Rates are Monte Carlo means, so exact ties almost never occur: candidates within $0.005$ of the best rate count as tied, and ties among equally worst communities go to the one with the most exposed nodes.
\end{itemize}

\subsection{Predictions from the ceiling model}
Claim 1 gives maximin's PoF $w_H+q\,w_L$ for $k\le E_H/b$ and $k\le k^*$.
\begin{itemize}
\item \textbf{Claim 3 (no leveling down).} Past this range, lexicographic maximin's PoF is non-increasing in $k$ and reaches $1$ at $k=E/b$, where both maximin and the optimum convert every acceptor ($a_H E_H+a_L E_L$). Throughout, the optimum's total $T^*$ is at least maximin's $T_m$, so the ratio $T^*/T_m$ cannot rise while $T^*$ grows no faster than $T_m$.
  \begin{itemize}
  \item For $E_H/b<k\le k^*$, the optimum has saturated $H$ and gains only $b\,a_L$ per seed, while maximin gains $bE/\sum_g(E_g/a_g)\ge b\,a_L$.
  \item Past $k^*$, $L$ is saturated at $r_L=a_L$, so the minimum is fixed and the tie-break sends every extra seed to $H$, where it saves $b\,a_H$, at least as much as an extra seed saves under the optimum.
  \item Plain maximin is indifferent among all allocations with $r_H\ge a_L$, so its worst case keeps $r_H=a_L$ and grows to DP's $q\,w_H+w_L$.
  \item So the tie-break, not the minimum, is what separates maximin from DP at large budgets.
  \end{itemize}
\item \textbf{Claim 4 (saved vs.\ converted).} By the converted bound (Section~\ref{sec:setup}), $r_c^{\mathrm{conv}}\le\bar a_c+\mathbb{E}\,|S_P\cap\text{exposed}_c|/E_c$, so $\mathbb{M}^{\mathrm{conv}}\le\min_c(\bar a_c+s_c)$, where $s_c=\mathbb{E}\,|S_P\cap\text{exposed}_c|/E_c$ is $c$'s seed share.
  \begin{itemize}
  \item The least receptive community caps the objective whatever the budget, and maximin over converted nodes pours seeds into it, where past $\bar a_c$ each seed adds little beyond itself.
  \item No such cap holds for $\mathbb{M}^{\mathrm{sav}}$, which can rise through shielding, so the two objectives can choose different worst communities, one set by acceptance and the other by network position (not observed with two groups).
  \end{itemize}
\end{itemize}

\subsection{Fragility and a soft minimum}
If more than $k$ infected communities have rates that only a seed inside them can raise, every $S_P\in\mathcal{S}$ has $\mathbb{M}=0$ and the first stage ranks nothing. With thousands of Louvain clusters (as in FWRRS), one tiny community no seed can help does the same.

We therefore also use a soft minimum
\[
\mathbb{M}_\tau=-\tau\log\sum_{c\in\mathcal{C}_{\mathrm{infected}}}e^{-r_c/\tau},
\]
which tends to $\mathbb{M}$ as $\tau\to0$ and gives greedy a nonzero gain from every seed.

\subsection{Experiments}
On the shared setup (Section~\ref{sec:setup}) we compare worst-first greedy, soft-min greedy ($\tau=0.05$), total greedy, DP greedy, the analytic allocation $k_g\propto E_g/\bar a_g$ and proportional seeding, on both outcomes (Figs.~\ref{fig:maximin-accept}--\ref{fig:maximin-fragility}).

\textbf{(A) Acceptance sweep.}
\begin{itemize}
\item \emph{Saved:} worst-first greedy keeps (almost) all its seeds in $H$, and its PoF is $0.97$--$0.99$ at every $\mu_L$, against the formula's $1.01$--$1.91$; values below $1$ reflect noise or total greedy's shortfall as a heuristic. Seeds in $H$ also shield $L$, so maximin is nearly free.
\item \emph{Converted:} worst-first's PoF ($1.09$--$1.81$) exceeds the formula for $\mu_L\ge0.45$ (e.g.\ $1.28$ vs.\ $1.16$), matches it at $\mu_L=0.3$ ($1.35$ vs.\ $1.35$) and falls below it at $\mu_L=0.15$ ($1.81$ vs.\ $1.91$). The formula's assumptions fail in this setup (Section~\ref{sec:parity-exp}), so these comparisons are only indicative. It stays well below DP greedy's at $\lambda=2$ ($2.47$ at $\mu_L=0.3$; Section~\ref{sec:compare} compares looser $\lambda$).
\item Worst-first is not the maximin optimum: soft-min greedy reaches a higher $\mathbb{M}^{\mathrm{conv}}$ at $\mu_L=0.3$ ($0.155$ vs.\ $0.137$), at a higher PoF ($1.69$). DP greedy at $\lambda=2$, with all ten seeds in $L$, is within noise of worst-first ($0.140$).
\end{itemize}

\textbf{(B) Budget sweep.}
\begin{itemize}
\item Worst-first's converted PoF is roughly flat ($1.35$--$1.37$) from $k=9$ to $20$, then falls to $1.29$ at $k=30$, consistent with Claim 3. Plain maximin with random ties, not the worst case Claim 3 bounds, pays $1.66\to1.40$ from $k=10$ to $30$, and DP at $\lambda=2$ $2.47\to1.46$.
\item At $k=30$ worst-first's $\mathbb{M}^{\mathrm{conv}}$ ($0.344$) is within about $0.01$ of DP's ($0.355$) and plain maximin's ($0.358$), but it saves $13\%$ more in total than DP and $9\%$ more than plain maximin, the tie-break effect Claim 3 predicts.
\item $L$'s converted rate passes $\bar a_L=0.28$ only through the seeds (27 of 30 in $L$), as Claim 4 allows.
\end{itemize}

\textbf{(C) Many communities.}
\begin{itemize}
\item No infected community was ever left at exactly $0$: rates are means over realizations, and the truth reaches every community a little.
\item The fragility shows up instead in worst-first greedy, which chases one small community at a time. With 124 Louvain communities its $\mathbb{M}^{\mathrm{sav}}$ ($0.37$) falls below fairness-blind total greedy's ($0.49$), and its total drops from $0.69$ to $0.50$.
\item Soft-min greedy holds up best ($\mathbb{M}^{\mathrm{sav}}=0.54$, total $0.69$). On the converted outcome every method's $\mathbb{M}$ falls to about $0.02$--$0.05$.
\item With many communities, $\mathbb{M}_\tau$ is a different objective. Since $\mathbb{M}-\tau\log|\mathcal{C}|\le\mathbb{M}_\tau\le\mathbb{M}$, the gap can reach $0.035$ with $2$ groups but $0.24$ with $124$ communities, more than the converted rates themselves. There $\mathbb{M}_\tau$ acts like an unweighted average over the communities near the bottom, regardless of their size, rather than a minimum.
\end{itemize}
